%========= INICIO DEL CÓDIGO PARA impacto_ambiental.tex =========


% -------------------------------------------------
% VERSION TRACKER — No modificar este bloque
% Versión:    Tesis Humanizada
% Archivo:    Capitulo7ImpactoAmbiental/impactoambiental.tex
% Última mod:  2026-05-08T08:23:24
% Git commit:  cc9c766f @ main
% Audiencia:   general
% Verificar antes de editar: bash check_tesis_version.sh overleaf_tesis_humanizado/Capitulo7ImpactoAmbiental/impactoambiental.tex
% -------------------------------------------------


El principal impacto ambiental de este proyecto es \textbf{positivo y se deriva de la desmaterialización}. Desmaterializar significa reemplazar objetos físicos ---papel, tinta, archivadores, cuadernos--- por equivalentes digitales. Al digitalizar la gestión de inventarios, el sistema reduce drásticamente la necesidad de registros en papel, impresiones de reportes y archivadores físicos.

Como orden de magnitud, una bodega podría dejar de imprimir algunas resmas de papel al año; si muchas PYMEs adoptaran una herramienta similar, el efecto acumulado sería mayor.

El impacto negativo es mínimo. Proviene del consumo eléctrico del computador usado durante el desarrollo y del servidor donde se despliega el sistema. Se estimó, como orden de magnitud y sin medición, un consumo de 600 kWh durante todo el proyecto, equivalente a unas 0.24 toneladas de dióxido de carbono (CO2). Este impacto se mitiga usando software eficiente y optimizando el código para que el servidor trabaje menos tiempo.

La estimación de la huella de carbono se calcula con el factor de conversión oficial del Sistema Interconectado Nacional de Colombia para 2021, establecido en 0.403 toneladas de CO2 equivalente por cada megavatio-hora (MWh) consumido \cite{upme_factor_emision_2021}. El cálculo es directo:
$$ 600\,\text{kWh} \times \frac{0.403\,\text{t CO2eq}}{1000\,\text{kWh}} \approx 0.24\,\text{toneladas de CO2eq.} $$

A futuro, este impacto podría reducirse aún más eligiendo un proveedor de hosting que use energías renovables.

%========= FIN DEL CÓDIGO PARA impacto_ambiental.tex =========
